\chapter{Números aleatorios}



En este capítulo se desarrollarán los siguientes temas:
\section*{Contenido}
\begin{enumerate}

\item Introducción.
\item Método Congruencial Binario.
\item Método de los Cuadrados Medios.
\item Método Congruencial Mixto.
\item Escogencia de Parámetros.
\item Método Congruencial Multiplicativo.
\item Método Congruencial Aditivo.
\item Método Congruencial Polinomial.
\item Generadores en Uso.
\item Resumen.
\item Ejercicios.

\end{enumerate}


\section{Introducción}

Un conjunto de números se denominan aleatorios cuando cada uno tiene la misma probabilidad de ser seleccionado. Esto significa que pertenecen a una distribución uniforme de probabilidad.

Es usual que se requieran secuencias de números aleatorios en un proceso de simulación. Para producir estos números existen tres posibles formas.

\begin{enumerate}
 \item Utilizar una provisión externa.

Se debe contar con una tabla o base de datos de números aleatorios. Esta debe estar almacenada en algún medio electrónico para accesarla. En los apéndices de esta obra se puede encontrar una tabla de números aleatorios que se utiliza para el planteamiento de algunos ejercicios.


\item Generación Interna por un proceso físico al azar.

Existen algunos dispositivos de hardware que son capaces de producir secuencias de números aleatorios. Al conectar estos dispositivos con las computadoras se logran producir programas cuyos eventos probabilísticos son imposibles de predecir.


\item Generación Interna por medio de algoritmos seudoaleatorios.

No existe forma de producir secuencias de números verdaderamente aleatorios por medio de un algoritmo computacional. Ya que una vez que se ha escrito el programa se puede deducir cuáles son los números que va a producir. 

Lo que se busca es escribir programas que produzcan secuencias de números que tengan las mismas propiedades que tendrían los números aleatorios. Este tipo de números suelen llamarse números seudoaleatorios, ya que son producto de un algoritmo computacional, pero pueden ser una aproximación útil a una secuencia verdaderamente aleatoria. En general la aproximación es útil cuando algunas propiedades importantes se conservan, mientras que otras propiedades pueden desestimarse.

\end{enumerate}


 

Sea cual sea el método que se utilice resulta útil que el generador de números seudoaleatorios cumpla con las siguientes características.
\begin{itemize}
\item Uniformemente distribuidos.
\item Estadísticamente independientes.
\item Reproducibles.
\item Baja repetición en sucesiones largas.
\item Generados por un método rápido.
\item Generados sin un alto nivel de almacenamiento.
\end{itemize}



\section{Método Congruencial Binario}


Este método para la generación de números aleatorios se basa en números binarios a los cuales se les aplican las operaciones de "xor" y "shift". La idea es iniciar con una secuencia cualquiera de números binarios, que no sean todos cero, y luego moverlos hacia la derecha una posición. La posición inicial se llena entonces con alguna operación.

Por ejemplo suponga que se inicia con la secuencia

$$1111$$

Si se hace un xor de los dos bits menos significativos, o sea las posiciones 0 y 1 de derecha a izquierda. Por lo tanto 1 xor 1 = 0, se genera un cero. Este valor se coloca como el $bit$ más significativo y se corren los demás una posición.

$$0111$$

Realicemos este proceso de forma sucesiva, a continuación se presentan los valores de $x(n)$ en base binaria y en base decimal. Observe que en este caso se obtienen todos los números entre 1 y 15\footnote{Es imposible generar el número 0, pues si todos los bits se convierten a cero se produce una parálisis del proceso de generación}.

\begin{verbatim}

---------------------
x( n)      B2     B10
---------------------
x( 0)  =  1111  =  15
x( 1)  =  0111  =   7
x( 2)  =  0011  =   3
x( 3)  =  0001  =   1
x( 4)  =  1000  =   8
x( 5)  =  0100  =   4
x( 6)  =  0010  =   2
x( 7)  =  1001  =   9
x( 8)  =  1100  =  12 
x( 9)  =  0110  =   6
x(10)  =  1011  =  11 
x(11)  =  0101  =   5
x(12)  =  1010  =  10
x(13)  =  1101  =  13
x(14)  =  1110  =  14
x(15)  =  1111  =  15
---------------------

\end{verbatim}

\section{Método de los Cuadrados Medios}

El método de los cuadrados medios consiste en el siguiente proceso:

\begin{enumerate}
\item Iniciar con una semilla entera.
\item Elevar el número al cuadrado.
\item Tomar una parte central de k dígitos.
\item Los k dígitos pasan a ser el nuevo número.
\end{enumerate}

Por ejemplo, si se desea calcular una secuencia con $k=4$ y $x(0) = 445$  se tiene el siguiente desarrollo.

\begin{verbatim}
-------------------------------------------------
x( 0)  =                                       445
x( 1)  =  (445)^2  =   198025  =  1-9802-5  = 9802
x( 2)  = (9802)^2  = 96079204  = 96-0792-04 = 0792
x( 3)  =  (792)^2  =   627264  =  6-2726-4  = 2726
x( 4)  = (2726)^2  =  7431076  =  7-4310-76 = 4310
x( 5)  = (4310)^2  = 18576100  = 18-5761-00 = 5761
x( 6)  = (5761)^2  = 33189121  = 33-1891-21 = 1891
x( 7)  = (1891)^2  =  3575881  =  3-5758-81 = 5758
x( 8)  = (5758)^2  = 33154564  = 33-1545-64 = 1545
x( 9)  = (1545)^2  =  2387025  =  2-3870-25 = 3870
x(10)  = (3870)^2  = 14976900  = 14-9769-00 = 9769
--------------------------------------------------
\end{verbatim}



\section{Método Congruencial Mixto}

El método congruencial mixto es una relación de recurrencia que inicia con una semilla. La semilla es un valor constante y luego se genera una sucesión matemática de la forma:

$$x(0)   = s$$ 
$$x(n+1) = f( x(n) )$$


Usualmente se establece la siguiente relación de recurrencia:

$$x(0)   = semilla$$
$$x(n+1) = ( a * x(n) + c ) (m)$$

Donde:

$x(0)$ : se denomina semilla\\
$a$    : se denomina multiplicador\\ 
$c$    : se denomina constante aditiva\\
$m$    : secuencia de números va desde $0$ hasta $m-1$\\

En los ejemplos que se muestran a continuación se utilizará la abreviatura $mod$, para referirse a la operación módulo. Por ejemplo, construyamos la secuencia que produciría la siguiente fórmula:
\begin{verbatim}

x(0)   = 4
x(n+1) = ( 5 * x(n) + 7 ) mod 8



x( 0)  =                                    4
x( 1)  =  (5*4 + 7) mod  8  =  27 mod 8  =  3
x( 2)  =  (5*3 + 7) mod  8  =  22 mod 8  =  6
x( 3)  =  (5*6 + 7) mod  8  =  37 mod 8  =  5
x( 4)  =  (5*5 + 7) mod  8  =  32 mod 8  =  0
x( 5)  =  (5*0 + 7) mod  8  =   7 mod 8  =  7
x( 6)  =  (5*7 + 7) mod  8  =  42 mod 8  =  2
x( 7)  =  (5*2 + 7) mod  8  =  17 mod 8  =  1
x( 8)  =  (5*1 + 7) mod  8  =  12 mod 8  =  4
x( 9)  =  (5*4 + 7) mod  8  =  27 mod 8  =  3
x(10)  =  (5*3 + 7) mod  8  =  22 mod 8  =  6
x(11)  =  (5*6 + 7) mod  8  =  37 mod 8  =  5
\end{verbatim}


Observe ahora lo que ocurre al seleccionar otros parámetros. En esta ocasión se produce un mal período ya que no se generan todos los números entre 0 y 9.
\begin{verbatim}
x(0)   = 7
x(n+1) = ( 7 * x(n) + 7 ) mod 10

--------------------------------------------
x(0)                                       7
x(1)  =  (7*7 + 7) mod 10  = 56 mod 10  =  6
x(2)  =  (7*6 + 7) mod 10  = 49 mod 10  =  9
x(3)  =  (7*9 + 7) mod 10  = 70 mod 10  =  0
x(4)  =  (7*0 + 7) mod 10  =  7 mod 10  =  7
--------------------------------------------
\end{verbatim}

Otro ejemplo de un mal período se presenta a continuación.

\begin{verbatim}
x(0)   = 0
x(n+1) = ( 19 * x(n) + 1) mod 381


---------------------------------------------------
x(0)  =                                           0
x(1)  =  (19* 0 + 1) mod 381  =    1 mod 381  =   1
x(2)  =  (19* 1 + 1) mod 381  =   20 mod 381  =  20
x(3)  =  (19*20 + 1) mod 381  =  381 mod 381  =   0
x(4)  =  (19* 0 + 1) mod 381  =    1 mod 381  =   1
x(5)  =  (19* 1 + 1) mod 381  =   20 mod 381  =  20
---------------------------------------------------
\end{verbatim}

Construyamos la secuencia que produciría la siguiente fórmula. Observe que en este caso se tiene todo el recorrido.

\begin{verbatim}
x(0)   = 5
x(n+1) = ( 1 * x(n) + 9 ) mod 10

----------------------------------------------
x( 0)  =                                     5
x( 1)  =  (1*5 + 9) mod 10  =  14 mod 10  =  4
x( 2)  =  (1*4 + 9) mod 10  =  13 mod 10  =  3
x( 3)  =  (1*3 + 9) mod 10  =  12 mod 10  =  2
x( 4)  =  (1*2 + 9) mod 10  =  11 mod 10  =  1
x( 5)  =  (1*1 + 9) mod 10  =  10 mod 10  =  0
x( 6)  =  (1*0 + 9) mod 10  =   9 mod 10  =  9
x( 7)  =  (1*9 + 9) mod 10  =  18 mod 10  =  8
x( 8)  =  (1*8 + 9) mod 10) =  17 mod 10  =  7
x( 9)  =  (1*7 + 9) mod 10) =  16 mod 10  =  6
x(10)  =  (1*6 + 9) mod 10) =  15 mod 10  =  5
----------------------------------------------
\end{verbatim}


Ahora en este nuevo ejemplo en apariencia se tiene un buen recorrido, sin embargo si nota el último dígito se tiene una secuencia similar a la anterior. Lo cual podría ser un problema dependiendo de la aplicación.


\begin{verbatim}
x(0)   = 5
x(n+1) = ( 81 * x(n) + 89 ) mod 100


------------------------------------------------------
x( 0)  =                                             5 
x( 1)  =  (81* 5 + 89) mod 100  =   494 mod 100  =  94
x( 2)  =  (81*94 + 89) mod 100  =  7703 mod 100  =  03
x( 3)  =  (81* 3 + 89) mod 100  =   332 mod 100  =  32
x( 4)  =  (81*32 + 89) mod 100  =  2681 mod 100  =  81
x( 5)  =  (81*81 + 89) mod 100  =  6650 mod 100  =  50
x( 6)  =  (81*50 + 89) mod 100  =  4139 mod 100  =  39
x( 7)  =  (81*39 + 89) mod 100  =  3248 mod 100  =  48
x( 8)  =  (81*48 + 89) mod 100  =  3977 mod 100  =  77
x( 9)  =  (81*77 + 89) mod 100  =  6326 mod 100  =  26
x(10)  =  (81*26 + 89) mod 100  =  2195 mod 100  =  95
x(11)  =  (81*95 + 89) mod 100  =  7784 mod 100  =  84
x(12)  =  (81*84 + 89) mod 100  =  6893 mod 100  =  93
x(13)  =  (81*93 + 89) mod 100  =  7622 mod 100  =  22
x(14)  =  (81*22 + 89) mod 100  =  1871 mod 100  =  71
x(15)  =  (81*71 + 89) mod 100  =  5840 mod 100  =  40
x(16)  =  (81*40 + 89) mod 100  =  3329 mod 100  =  29
x(17)  =  (81*29 + 89) mod 100  =  2438 mod 100  =  38
x(18)  =  (81*38 + 89) mod 100  =  3167 mod 100  =  67
x(19)  =  (81*67 + 89) mod 100  =  5516 mod 100  =  16
x(20)  =  (81*16 + 89) mod 100  =  1385 mod 100  =  85
------------------------------------------------------
\end{verbatim}

De los ejemplos anteriores se advierte el cuidado que se debe tener en el momento de escoger los parámetros del método para que el generador tenga un período completo. Este método debe considerar la coprimalidad de las partes entre otras cosas.




\section{Escogencia de Parámetros}


A pesar de la sencillez de este método, ha sido el objeto de gran cantidad de análisis matemático para determinar qué valores de los parámetros son apropiados.


\subsection{ Consideraciones Generales}



Para que el método congruencial mixto funcione adecuadamente se deben tener una serie de restricciones elementales.
\begin{itemize}
\item m > x(0)
\item m > a
\item m > c
\end{itemize}


\subsection{ Escogencia de m:}

\begin{enumerate}

\item El valor de m debe ser lo más grande posible, preferente del tamaño de palabra del computador. Pero puede ser más pequeño si causa problemas numéricos o de $overflow$.

\item También se han obtenido buenos resultados poniendo m como el número primo más grande posible tal que:

$$m < p^d$$

donde p es la base del sistema, en este caso se trata de un sistema decimal y d es el número de bits que tiene la palabra del computador.

\item Otros buenos resultados se presentan cuando se utiliza como m como valores de potencias de 2 o de 10.

\item Como un caso particular del punto anterior se han encontrado algunos valores buenos de m como:

$$m = 2^{31}-1 =  2,147,483,647$$

$$m = 10^{8} =  100000000$$


\item En algunos casos el valor de m está determinado por el problema. Por ejemplo si se desea simular el comportamiento de un dado se puede pensar que m debe ser 6, o un múltiplo de 6.

\end{enumerate}

\subsection{Selección del valor de a}

\begin{enumerate}
\item El valor de a no debe ser ni muy grande ni muy pequeño.

\item En algunos casos, si m es escogido adecuadamente, a puede tomar el valor de m-1. (coprimo con m)

\item El valor de a debe ser un entero impar y no debe ser divisible por 3 o 5.

\item Como un caso particular y muy peculiar se han encontrado buenos resultados con números que terminen en 21. Por ejemplo 821, 7021, etc.
\end{enumerate}

\subsection{Selección de c}

\begin{enumerate}
\item En general c puede ser cualquier valor constante, se ha encontrado que el valor de 1 es bueno.

\item En base binaria c = 5 es bueno y en base decimal c = 21 trabajan bien.
\end{enumerate}


\subsection{Selección de la $x_0$}
-----------------------

En general no es relevante, salvo la restricción mencionada que $m > x(0)$.



\subsection{Un buen generador general.}

Se ha demostrado que un buen generador que se usa a menudo tiene los siguientes parámetros:

$$a = 7^5  = 16,807$$

$$c = 0$$

$$m = 2^{31}-1 = 2147483647$$



Los criterios anteriores se pueden resumir en la siguiente tabla, que se toma de [Brysbaert,1991], muestra algunos generadores aleatorios que producen buenos resultados.

\begin{verbatim}
-------------------------------------------------------
x(0)                a       c      m   Autor
-------------------------------------------------------
impar           13^13       0   2^59    Ripley(1983)
                              
impar  44485709377909       0   2^48    Ripley(1983)
                              
impar            5^15       0   2^35    Knuth(1981)
                              
impar           69069   impar   2^32    Ripley(1983)
                              
impar          100485   impar   2^32    Atkinson(1980)
                              
impar      1589013525       1   2^32    Ripley(1983)
                              
impar       663608941   impar   2^32    Ripley(1983)
                              
<> 0        630360016       0   2^31-1  Dudewicz(1981)
                              
<> 0        764261123       0   2^31-1  Dudewicz(1981)
                              
<> 0        950706376       0   2^31-1  Fisherman(1986)
-------------------------------------------------------
\end{verbatim}



\section{Método Congruencial Multiplicativo}

Una vez que se ha presentado el método congruencial mixto se pueden realizar una serie de variantes. Por ejemplo el método congruencial multiplicativo en el cual c = 0.

$$x(0)   = semilla$$
$$x(n+1) = (a * x(n)) mod (m)$$


Por ejemplo:
\begin{verbatim}
x(0)   = 47
x(n+1) = (13 * x(n)) mod 767

-----------------------------------------------------
x( 0)                                          =   47
x( 1)  =  (13 *  47) mod 767  =   611 mod 767  =  611
x( 2)  =  (13 * 611) mod 767  =  7943 mod 767  =  273
x( 3)  =  (13 * 273) mod 767  =  3549 mod 767  =  481
x( 4)  =  (13 * 481) mod 767  =  6253 mod 767  =  117
x( 5)  =  (13 * 117) mod 767  =  1521 mod 767  =  754
x( 6)  =  (13 * 754) mod 767  =  9802 mod 767  =  598
x( 7)  =  (13 * 598) mod 767  =  7774 mod 767  =  104
x( 8)  =  (13 * 104) mod 767  =  1352 mod 767  =  585
x( 9)  =  (13 * 585) mod 767  =  7605 mod 767  =  702
x(10)  =  (13 * 702) mod 767  =  9126 mod 767  =  689
-----------------------------------------------------
\end{verbatim}



\section{Método Congruencial Aditivo}

Otra variante es el método congruencial aditivo. Esta técnica se basa en la fórmula de recurrencia dada por:

$$x(0)   = semilla_0$$
$$x(1)   = semilla_1$$
$$x(n+1) = (x(n) + x(n-1)) mod (m)$$


Observe que si x(0) y x(1) se inicializan con 1 se produce la conocida serie de números de Fibonacci.

\begin{verbatim}
x(0)   = 1
x(1)   = 1
x(n+1) = (x(n) + x(n-1)) mod 231


------------------------------------------------------
x( 0)                                           =    1
x( 1)  =                                        =    1
x( 2)  =  (  1 +   1) mod 231  =    2  mod 231  =    2
x( 3)  =  (  2 +   1) mod 231  =    3  mod 231  =    3
x( 4)  =  (  3 +   2) mod 231  =    5  mod 231  =    5
x( 5)  =  (  5 +   3) mod 231  =    8  mod 231  =    8
x( 6)  =  (  8 +   5) mod 231  =   13  mod 231  =   13
x( 7)  =  ( 13 +   8) mod 231  =   21  mod 231  =   21
x( 8)  =  ( 21 +  13) mod 231  =   34  mod 231  =   34
x( 9)  =  ( 34 +  21) mod 231  =   55  mod 231  =   55
x(10)  =  ( 55 +  34) mod 231  =   89  mod 231  =   89
------------------------------------------------------
\end{verbatim}


Otro ejemplo puede estar dado por:

\begin{verbatim}
x(0)   = 5
x(1)   = 6
x(n+1) = (x(n) + x(n-1)) mod 231


------------------------------------------------------
x( 0)                                           =    5
x( 1)  =                                        =    6
x( 2)  =  (  6 +   5) mod 231  =   11  mod 231  =   11
x( 3)  =  ( 11 +   6) mod 231  =   17  mod 231  =   17
x( 4)  =  ( 17 +  11) mod 231  =   28  mod 231  =   28
x( 5)  =  ( 28 +  17) mod 231  =   45  mod 231  =   45
x( 6)  =  ( 45 +  28) mod 231  =   73  mod 231  =   73
x( 7)  =  ( 73 +  45) mod 231  =  118  mod 231  =  118
x( 8)  =  (118 +  73) mod 231  =  191  mod 231  =  191
x( 9)  =  (191 + 118) mod 231  =  309  mod 231  =   78
x(10)  =  ( 78 + 191) mod 231  =  269  mod 231  =   38
------------------------------------------------------
\end{verbatim}



\section{Método Congruencial Polinomial}


Uno de los métodos polinomiales más conocidos es el generador congruencial cuadrático. Esta técnica se basa en la fórmula de recurrencia dada por:

$$x(0)   = semilla$$
$$x(n+1) = (d * x(n)^2 + a * x(n) + c) mod (m)$$


A continuación se presenta un ejemplo del método.
\begin{verbatim}
x(0)   = 0
x(n+1) = (51 * x(n)^2 + 7 * x(n) + 237) mod 100


-------------------------------------------------
x( 0)                                        =  0
x( 1)  = (51 *  0^2 +  7 *  0 + 237) mod 100 = 37
x( 2)  = (51 * 37^2 +  7 * 37 + 237) mod 100 = 15
x( 3)  = (51 * 15^2 +  7 * 15 + 237) mod 100 = 17
x( 4)  = (51 * 17^2 +  7 * 17 + 237) mod 100 = 95
x( 5)  = (51 * 95^2 +  7 * 91 + 237) mod 100 = 77
x( 6)  = (51 * 77^2 +  7 * 77 + 237) mod 100 = 55
x( 7)  = (51 * 55^2 +  7 * 55 + 237) mod 100 = 97
x( 8)  = (51 * 97^2 +  7 * 97 + 237) mod 100 = 75
x( 9)  = (51 * 75^2 +  7 * 75 + 237) mod 100 = 37
x(10)  = (51 * 37^2 +  7 * 37 + 237) mod 100 = 15
------------------------------------------------------
\end{verbatim}

Una de sus principales desventajas es el costo computacional, ya que son mucho más lentos que los métodos anteriores. Sin embargo suelen producir secuencias de números con muchos menos repeticiones y de períodos más largos.



\section{Generadores en Uso}

La biblioteca estándar de Unix utiliza el siguiente generador. En Unix, los números aleatorios se presentan en el intervalo [0,1] por lo que se divide el número por su recorrido para quedar en ese intervalo.

$$x(n+1) = (25214903917 \times x(n) + 11) mod (2^{48})$$

$$u(n+1) = \frac{x(n+1)}{2^{48}}$$


Otros lenguajes de programación utilizan el generador:

$$x(n+1) = (1140671485 \times x(n) + 12820163) mod (2^{24})$$

$$u(n+1) = \frac{x(n+1)}{2^{24}}$$


Además existe un website del cual se pueden obtener números aleatorios reales, no seudoaleatorios. Tal como se indicó sería un mecanismo de provisión externa de números. Este website se encuentra en:
\textit{www.random.org}


\section{Resumen}


En el presente capítulo se han presentado varios métodos para la generación de números aleatorios, el método congruencial binario, el método de los cuadrados medios, el método congruencial mixto y algunas de sus variantes.

Cada método tiene ciertos parámetros con los que funciona adecuadamente o no. En los siguientes capítulos se analizarán procesos estadísticos para evaluar el desempeño de estos generadores aleatorios.



 
\section{Ejercicios}.

\begin{enumerate}
    \renewcommand{\theenumi}{Ejercicio \arabic{enumi}} % Cambia el formato de numeración
    \setcounter{enumi}{0} % Inicia la numeración desde 1
    \item Utilice el método congruencial binario pero en lugar de realizar un xor con las posiciones 0 y 1, de derecha a izquierda, realice un xor con las posiciones 0 y 2. Utilice las siguientes semillas y muestre si realiza todo el recorrido o no.
    \begin{enumerate}[label=\alph*.]
    \item $x_0 = 110$
    \item $x_0 = 1111$
    \end{enumerate}
    \item Para 5 bits construya un método congruencial binario que realice un recorrido completo. Indique cuál es la semilla y cuál es la operación que realiza.

    \item Genere 50 números aleatorios de 4 dígitos utilizando el generador de cuadrados medios utilizando los siguientes parámetros.
    \begin{enumerate}[label=\alph*.]
    \item $x_0 = 3567345$
    \item $x_0 = 1234500012$
    \item $x_0 = 4567234902$
    \end{enumerate}
    \item Determine para cada generador congruencial mixto que se presenta a continuación si cumple con todo el período de m.
    \begin{enumerate}[label=\alph*.]
    \item $x_0 = 7$ \\$x_{n+1} = (5 \times x_n + 24)(32)$
    \item $x_0 = 8$ \\$x_{n+1} = (9 \times x_n + 13)(32)$
    \item $x_0 = 13$ \\$x_{n+1} = (50 \times x_n + 17)(64)$
    \item $x_0 = 15$ \\$x_{n+1} = (8 \times x_n + 16)(100)$
    \item $x_0 = 3$ \\$x_{n+1} = (5 \times x_n + 21)(100)$
    \end{enumerate}
    \item Determine para cada generador congruencial multiplicativo que se presenta a continuación si cumple con todo el período de m. La expresión $a^b$ significa a elevado a la potencia b. 
    \begin{enumerate}[label=\alph*.]
    \item $x_0 = 7$\\$x_{n+1} = (5 \times x_n)(2^6)$
    \item $x_0 = 9$\\$x_{n+1} = (11 \times x_n)(2^7)$
    \item $x_0 = 3$\\$x_{n+1} = (221 \times x_n)(10^3)$
    \item $x_0 = 17$\\$x_{n+1} = (203 \times x_n)(10^5)$
    \item $x_0 = 19$\\$x_{n+1} = (211 \times x_n)(10^8)$
    \end{enumerate}
    \item Utilice el método congruencial mixto y encuentre los parámetros adecuados de x(0), b, c y m para generar un recorrido completo de números entre 0 y 14.
    \item Defina los parámetros de a, c, x(0) para un recorrido completo con el valor de m = 9.
    \item Defina los parámetros de a, c, x(0) para un recorrido completo con el valor de m = 19.
    \item Utilizando las variantes del método congruencial mixto, el método congruencial multiplicativo y el método congruencial aditivo encuentre generadores para números de 0 hasta 9 que tengan un recorrido completo.
    \item Se tiene un generador aleatorio de un sólo dígito. Es decir puede generar de manera uniforme los valores $0,1,2,3,4,5,6,7,8$ y $9$. Para cada uno de los siguientes fenómenos indique como utilizaría este generador.
    \begin{enumerate}[label=\alph*.]
    \item Para generar una moneda legal, con 0.50 de probabilidad de que salga corona y 0.50 de probabilidad que salga escudo.
    \item Para generar el siguiente fenómeno, probabilidad de a es 0.10, la probabilidad de b es 0.40 y la probabilidad de c es 0.50.
    \item Para generar los valores de arrojar un dado (uniforme de seis caras)
    \item Si en lugar de un generador aleatorio entre $[0,9]$ se tiene un generador de números aleatorios entre $[0,99]$ indique cómo realizaría la asignación de probabilidades anterior y considere ¿Se pueden simular dados con todos los poliedros regulares?.
    \end{enumerate}
\end{enumerate}






